\chapter{Diagonalization}

In the last chapter you learned to find a matrix's eigenvectors: the special
directions a machine does not turn, only stretch. An eigenvector $\vec{v}$ with
eigenvalue $\lambda$ satisfies the one equation that this whole chapter rests on,
\[
  \mat{A}\vec{v} = \lambda\,\vec{v}.
\]
That is one good fact about \emph{one} direction. The follow-up question is the
one that pays off: \emph{what if we use the eigenvectors as our coordinate axes?}
Instead of describing space with the standard grid $\vec{e}_1, \vec{e}_2, \dots$,
we lay down a new grid built entirely from the directions $\mat{A}$ likes.

In that grid the matrix $\mat{A}$, which had mixed up the coordinates and looked
complicated, does nothing but scale each axis by its own factor. The tangled
machine becomes a set of independent dials. This change into the eigenbasis is
called \textbf{diagonalization}, and it is the workhorse of applied linear
algebra. It computes $\mat{A}^{100}$ without a hundred multiplications, solves
recurrences like Fibonacci in closed form, and explains systems of differential
equations; a few chapters from now it becomes the seed of the singular value
decomposition that runs modern data science.

We build the picture geometrically first, then write down the formula
$\mat{A} = \mat{S}\mat{\Lambda}\mat{S}\inv$ and push it hard: powers, the
geometry of working in the eigenbasis, the matrices that \emph{refuse} to
diagonalize, and a closed-form Fibonacci you can check on your fingers.

\section{The Big Idea: Work in the Eigenbasis}
\label{ch24:big-idea}

Start with the picture before the formula. Suppose a $2\times 2$ matrix
$\mat{A}$ has two eigenvectors $\vec{v}_1$ and $\vec{v}_2$ pointing in different
directions, with eigenvalues $\lambda_1$ and $\lambda_2$. Because they point in
different directions, these two arrows are a \emph{basis}: every vector $\vec{x}$
in the plane can be written as a combination of them,
\[
  \vec{x} = c_1\vec{v}_1 + c_2\vec{v}_2.
\]
Now feed $\vec{x}$ to the machine. Linearity does all the work:
\[
  \mat{A}\vec{x}
  = \mat{A}(c_1\vec{v}_1 + c_2\vec{v}_2)
  = c_1\,\mat{A}\vec{v}_1 + c_2\,\mat{A}\vec{v}_2
  = c_1\lambda_1\vec{v}_1 + c_2\lambda_2\vec{v}_2.
\]
In the eigenvector coordinates $(c_1, c_2)$, the matrix did nothing but multiply
the first coordinate by $\lambda_1$ and the second by $\lambda_2$. No mixing. The
complicated machine, viewed in the right coordinates, is just two independent
stretches.

\begin{intuitionbox}
A general matrix mixes your coordinates together: the new $x$ depends on the old
$x$ \emph{and} the old $y$. Its eigenvectors are the directions where this mixing
stops. If you rebuild space out of eigenvectors -- use them as your axes -- then
$\mat{A}$ stops being a tangle and becomes a row of separate volume knobs, one
per axis, each turned to its eigenvalue $\lambda_i$. \textbf{Diagonalization is
the act of choosing the coordinate system in which the machine is simplest.}
\end{intuitionbox}

\begin{figure}[H]
\centering
\begin{tikzpicture}[scale=1.0,>=Stealth]
  % Left: standard grid getting sheared/stretched in a confusing way
  \begin{scope}[shift={(0,0)}]
    \draw[->,gray!55] (-0.3,0) -- (3.4,0) node[right]{};
    \draw[->,gray!55] (0,-0.3) -- (0,3.4) node[above]{};
    % a generic vector and its (mixed-up) image under A
    \draw[->,vecblue,very thick] (0,0) -- (1.6,0.8) node[above right]{$\vec{x}$};
    \draw[->,vecblue!55,dashed,very thick] (0,0) -- (2.6,2.2) node[right]{$\mat{A}\vec{x}$};
    \node[below] at (1.6,-0.3) {\small standard view: $\mat{A}$ mixes axes};
  \end{scope}
  % Right: eigen grid, A just scales along each axis
  \begin{scope}[shift={(6,0)}]
    \draw[->,vecblue!50,thin] (0,0) -- (3.2,1.6);
    \draw[->,vecred!50,thin] (0,0) -- (1.28,3.2);
    \draw[->,vecblue,very thick] (0,0) -- (1.6,0.8) node[midway,below right]{$\vec{v}_1$};
    \draw[->,vecred,very thick] (0,0) -- (0.8,2.0) node[midway,left]{$\vec{v}_2$};
    % scaled images along each eigen-axis (dashed)
    \draw[->,vecblue,dashed,thick] (1.6,0.8) -- (3.2,1.6) node[right]{$\lambda_1\vec{v}_1$};
    \draw[->,vecred,dashed,thick] (0.8,2.0) -- (1.28,3.2) node[above]{$\lambda_2\vec{v}_2$};
    \node[below] at (1.7,-0.2) {\small eigen view: each axis just scales};
  \end{scope}
\end{tikzpicture}
\caption{Left: in the standard grid, $\mat{A}$ looks like it stirs the $x$ and
$y$ coordinates together -- a generic $\vec{x}$ is sent to a tilted, rescaled
$\mat{A}\vec{x}$ with no obvious pattern. Right: along the eigenvectors
$\vec{v}_1,\vec{v}_2$ the
action is pure scaling -- the $\vec{v}_1$ axis stretches by $\lambda_1$, the
$\vec{v}_2$ axis by $\lambda_2$. Same machine, smarter axes.}
\label{ch24:fig-eigenview}
\end{figure}

\begin{focusbox}
The slogan for the whole chapter: \textbf{eigenvectors are the axes in which a
matrix is diagonal.} Find those axes and the matrix loses all its complexity.
\end{focusbox}

\section{The Factorization \texorpdfstring{$\mat{A}=\mat{S}\mat{\Lambda}\mat{S}\inv$}{A = S Lambda S inverse}}
\label{ch24:factorization}

Now the algebra. The trick is to package all the eigenvectors into the columns of
a single matrix and watch what $\mat{A}$ does to them all at once.

Suppose $\mat{A}$ is $n\times n$ and has $n$ eigenvectors
$\vec{v}_1, \dots, \vec{v}_n$ with eigenvalues $\lambda_1, \dots, \lambda_n$.
Build the \textbf{eigenvector matrix}
\[
  \mat{S} = \big[\; \vec{v}_1 \;\big|\; \vec{v}_2 \;\big|\; \cdots \;\big|\; \vec{v}_n \;\big],
\]
putting one eigenvector in each column. Now multiply $\mat{A}$ by $\mat{S}$.
Matrix times matrix works column by column -- the $j$-th column of $\mat{A}\mat{S}$
is $\mat{A}$ times the $j$-th column of $\mat{S}$, which is $\mat{A}\vec{v}_j$.
But $\vec{v}_j$ is an eigenvector, so $\mat{A}\vec{v}_j = \lambda_j\vec{v}_j$.
Therefore
\[
  \mat{A}\mat{S}
  = \big[\; \lambda_1\vec{v}_1 \;\big|\; \lambda_2\vec{v}_2 \;\big|\; \cdots \;\big|\; \lambda_n\vec{v}_n \;\big].
\]
Each column got scaled by its own eigenvalue. That right-hand side is exactly
$\mat{S}$ with its columns rescaled, and you
rescale the columns of a matrix by multiplying \emph{on the right} by a diagonal
matrix. Let
\[
  \mat{\Lambda} = \diag(\lambda_1, \dots, \lambda_n)
  = \begin{bmatrix} \lambda_1 & & \\ & \ddots & \\ & & \lambda_n \end{bmatrix}.
\]
Then $\mat{S}\mat{\Lambda}$ takes column $j$ of $\mat{S}$ and multiplies it by
$\lambda_j$ -- precisely the right-hand side above. That proves the central
identity
\[
  \boxed{\;\mat{A}\mat{S} = \mat{S}\mat{\Lambda}.\;}
\]

\begin{trapbox}
Order matters. $\mat{S}\mat{\Lambda}$ (diagonal on the \emph{right}) scales the
\emph{columns} of $\mat{S}$. $\mat{\Lambda}\mat{S}$ (diagonal on the
\emph{left}) would scale the \emph{rows} -- a different matrix entirely. The
eigenvalue must multiply the eigenvector it belongs to, so it has to sit on the
right: $\mat{A}\mat{S} = \mat{S}\mat{\Lambda}$, not $\mat{\Lambda}\mat{S}$.
\end{trapbox}

So far we needed no inverse. The inverse enters only when we want to
\emph{solve} for $\mat{A}$. If the eigenvectors are independent, then $\mat{S}$ is
invertible (its columns are a basis), and we may multiply $\mat{A}\mat{S} =
\mat{S}\mat{\Lambda}$ on the right by $\mat{S}\inv$ to isolate $\mat{A}$, or on the
left by $\mat{S}\inv$ to isolate $\mat{\Lambda}$:
\[
  \boxed{\;\mat{A} = \mat{S}\,\mat{\Lambda}\,\mat{S}\inv\;}
  \qquad\text{and}\qquad
  \boxed{\;\mat{\Lambda} = \mat{S}\inv\,\mat{A}\,\mat{S}.\;}
\]

\begin{definition}[Diagonalizable matrix]
\label{ch24:def-diagonalizable}
An $n\times n$ matrix $\mat{A}$ is \textbf{diagonalizable} if it can be written
$\mat{A} = \mat{S}\mat{\Lambda}\mat{S}\inv$ for some invertible $\mat{S}$ and
diagonal $\mat{\Lambda}$. Equivalently, $\mat{A}$ is diagonalizable if and only
if it has $n$ linearly independent eigenvectors. Those eigenvectors are the
columns of $\mat{S}$; the eigenvalues are the diagonal entries of
$\mat{\Lambda}$, in the matching order.
\end{definition}

\begin{formalbox}
\textbf{Theorem (Diagonalization).} An $n\times n$ matrix $\mat{A}$ is
diagonalizable over $\R$ if and only if $\R^n$ has a basis consisting of
eigenvectors of $\mat{A}$. In that case, writing those basis vectors as the
columns of $\mat{S}$ and their eigenvalues on the diagonal of $\mat{\Lambda}$,
\[
  \mat{A} = \mat{S}\mat{\Lambda}\mat{S}\inv .
\]
The factorization is not unique: you may reorder the eigenpairs (permuting
columns of $\mat{S}$ and the matching diagonal entries of $\mat{\Lambda}$
together) or rescale any eigenvector, and you still get $\mat{A}$ back.
\end{formalbox}

\begin{strangbox}
Strang writes this factorization with the same letters: $\mat{A} =
\mat{S}\mat{\Lambda}\mat{S}\inv$, with $\mat{S}$ the eigenvector matrix and
$\mat{\Lambda}$ (capital lambda) the eigenvalue matrix. He stresses, as we will,
that ``diagonalizing requires $n$ independent eigenvectors.'' If you read this
\emph{before} Strang, the one thing to carry in: $\mat{S}$'s columns are
eigenvectors and $\mat{\Lambda}$'s diagonal is eigenvalues, \emph{in the same
order}. If you read it \emph{after}, notice we are deliberately separating two
hypotheses Strang sometimes bundles: ``has $n$ eigenvalues'' versus ``has $n$
independent eigenvectors.'' Only the second is what diagonalization needs.
\end{strangbox}

\subsection{A first worked example}
\label{ch24:first-example}

Take a concrete matrix from start to finish.
\[
  \mat{A} = \begin{bmatrix} 2 & 1 \\ 1 & 2 \end{bmatrix}.
\]
\textbf{Step 1: eigenvalues.} Solve $\det(\mat{A} - \lambda\mat{I}) = 0$:
\[
  \det\begin{bmatrix} 2-\lambda & 1 \\ 1 & 2-\lambda \end{bmatrix}
  = (2-\lambda)^2 - 1
  = \lambda^2 - 4\lambda + 3
  = (\lambda - 1)(\lambda - 3) = 0,
\]
so $\lambda_1 = 1$ and $\lambda_2 = 3$.

\textbf{Step 2: eigenvectors.} For $\lambda_1 = 1$, solve $(\mat{A}-\mat{I})\vec{v}=\vec0$:
\[
  \begin{bmatrix} 1 & 1 \\ 1 & 1 \end{bmatrix}\vec{v} = \vec0
  \;\Longrightarrow\; v_1 + v_2 = 0
  \;\Longrightarrow\; \vec{v}_1 = \colvec{1\\-1}.
\]
For $\lambda_2 = 3$, solve $(\mat{A}-3\mat{I})\vec{v}=\vec0$:
\[
  \begin{bmatrix} -1 & 1 \\ 1 & -1 \end{bmatrix}\vec{v} = \vec0
  \;\Longrightarrow\; -v_1 + v_2 = 0
  \;\Longrightarrow\; \vec{v}_2 = \colvec{1\\1}.
\]

\textbf{Step 3: assemble.} Place eigenvectors as columns, eigenvalues on the diagonal:
\[
  \mat{S} = \begin{bmatrix} 1 & 1 \\ -1 & 1 \end{bmatrix},
  \qquad
  \mat{\Lambda} = \begin{bmatrix} 1 & 0 \\ 0 & 3 \end{bmatrix}.
\]
The inverse of a $2\times2$ matrix $\begin{bsmallmatrix} a&b\\c&d\end{bsmallmatrix}$
is $\frac{1}{ad-bc}\begin{bsmallmatrix} d&-b\\-c&a\end{bsmallmatrix}$, and here
$ad-bc = (1)(1)-(1)(-1) = 2$, so
\[
  \mat{S}\inv = \frac{1}{2}\begin{bmatrix} 1 & -1 \\ 1 & 1 \end{bmatrix}.
\]

\textbf{Step 4: check.} Multiply $\mat{S}\mat{\Lambda}\mat{S}\inv$ back together
and confirm we recover $\mat{A}$:
\[
  \mat{S}\mat{\Lambda}
  = \begin{bmatrix} 1 & 1 \\ -1 & 1 \end{bmatrix}\begin{bmatrix} 1 & 0 \\ 0 & 3 \end{bmatrix}
  = \begin{bmatrix} 1 & 3 \\ -1 & 3 \end{bmatrix},
\]
\[
  \mat{S}\mat{\Lambda}\mat{S}\inv
  = \begin{bmatrix} 1 & 3 \\ -1 & 3 \end{bmatrix}\cdot\frac{1}{2}\begin{bmatrix} 1 & -1 \\ 1 & 1 \end{bmatrix}
  = \frac{1}{2}\begin{bmatrix} 4 & 2 \\ 2 & 4 \end{bmatrix}
  = \begin{bmatrix} 2 & 1 \\ 1 & 2 \end{bmatrix} = \mat{A}.\ \checkmark
\]

\begin{labbox}
Diagonalize $\mat{B} = \begin{bsmallmatrix} 3 & 0 \\ 1 & 2 \end{bsmallmatrix}$ by
hand. It is lower triangular, so read the eigenvalues straight off the diagonal:
$\lambda = 3$ and $\lambda = 2$. Find one eigenvector for each (you should get
$\colvec{1\\1}$ for $\lambda=3$ and $\colvec{0\\1}$ for $\lambda=2$), assemble
$\mat{S}$ and $\mat{\Lambda}$, and multiply $\mat{S}\mat{\Lambda}\mat{S}\inv$
back to verify. Doing this once by hand cements the recipe forever.
\end{labbox}

\section{Why Powers Become Easy}
\label{ch24:powers}

This is the payoff. Suppose you need $\mat{A}^k$ for large $k$ -- maybe you are
iterating a process a thousand times. Multiplying $\mat{A}$ by itself a thousand
times is slow and numerically dangerous. Diagonalization makes it a one-liner.

Watch the inverses cancel. Square the factorization:
\[
  \mat{A}^2
  = (\mat{S}\mat{\Lambda}\mat{S}\inv)(\mat{S}\mat{\Lambda}\mat{S}\inv)
  = \mat{S}\mat{\Lambda}\underbrace{(\mat{S}\inv\mat{S})}_{=\,\mat{I}}\mat{\Lambda}\mat{S}\inv
  = \mat{S}\mat{\Lambda}^2\mat{S}\inv.
\]
The interior $\mat{S}\inv\mat{S}$ collapses to the identity. Do it $k$ times and
every interior pair collapses, leaving
\[
  \boxed{\;\mat{A}^k = \mat{S}\,\mat{\Lambda}^k\,\mat{S}\inv.\;}
\]
And raising a \emph{diagonal} matrix to a power is trivial -- you just raise each
diagonal entry to that power, because a diagonal matrix times itself only
multiplies corresponding entries:
\[
  \mat{\Lambda}^k
  = \begin{bmatrix} \lambda_1^k & & \\ & \ddots & \\ & & \lambda_n^k \end{bmatrix}.
\]
No matrix multiplication chain at all -- just $n$ scalar powers, then two matrix
products to sandwich them back. The hard part (the eigenvectors) is computed once
and reused for every $k$.

\begin{intuitionbox}
Think of $\mat{S}\inv$ as ``translate into eigen-language,'' $\mat{\Lambda}^k$ as
``do the simple thing $k$ times -- just stretch each axis,'' and $\mat{S}$ as
``translate back to ordinary coordinates.'' Powers are hard in ordinary
coordinates and trivial in eigen-coordinates, so we step into that language, do
the easy thing, and step back out.
\end{intuitionbox}

\begin{example}[A power computed by hand]
\label{ch24:power-example}
Use our $\mat{A} = \begin{bsmallmatrix} 2&1\\1&2\end{bsmallmatrix}$ from above to
find $\mat{A}^{10}$ without ten multiplications. We have $\lambda_1=1$,
$\lambda_2=3$, so $\mat{\Lambda}^{10} = \begin{bsmallmatrix} 1^{10}&0\\0&3^{10}\end{bsmallmatrix}
= \begin{bsmallmatrix} 1 & 0 \\ 0 & 59049 \end{bsmallmatrix}$. Then
\[
  \mat{A}^{10}
  = \mat{S}\mat{\Lambda}^{10}\mat{S}\inv
  = \begin{bmatrix} 1 & 1 \\ -1 & 1 \end{bmatrix}
    \begin{bmatrix} 1 & 0 \\ 0 & 59049 \end{bmatrix}
    \frac{1}{2}\begin{bmatrix} 1 & -1 \\ 1 & 1 \end{bmatrix}.
\]
Carry it out: $\mat{S}\mat{\Lambda}^{10} = \begin{bsmallmatrix} 1 & 59049 \\ -1 & 59049 \end{bsmallmatrix}$, and
\[
  \mat{A}^{10}
  = \frac{1}{2}\begin{bmatrix} 1 & 59049 \\ -1 & 59049 \end{bmatrix}\begin{bmatrix} 1 & -1 \\ 1 & 1 \end{bmatrix}
  = \frac{1}{2}\begin{bmatrix} 59050 & 59048 \\ 59048 & 59050 \end{bmatrix}
  = \begin{bmatrix} 29525 & 29524 \\ 29524 & 29525 \end{bmatrix}.
\]
Notice the structure the eigenvalues predicted: the entries are huge (driven by
$3^{10}$) and the matrix is nearly $\tfrac{3^{10}}{2}$ times the rank-one
``all-ones'' pattern, because the $\lambda=3$ direction dominates. That dominance
is the seed of the power method and of the long-run behavior of Markov chains.
\end{example}

\begin{aestheticbox}
A power $\mat{A}^k$ looks like it should grow in complexity with $k$ -- more and
more multiplications, entries spiraling out of control. Diagonalization shows that
nothing complicated is happening at all: the matrix is secretly $n$ independent
numbers $\lambda_i$, each raised to the $k$. All the apparent complexity lived in
the \emph{choice of coordinates}, not in the matrix's actual behavior. Finding the
eigenbasis is finding the coordinates in which the behavior is simple.
\end{aestheticbox}

\subsection{The long-run behavior you can read off}
\label{ch24:longrun}

Because $\mat{A}^k$ scales the $\vec{v}_i$ direction by $\lambda_i^k$, the
\emph{largest} $\abs{\lambda_i}$ wins as $k\to\infty$. If one eigenvalue is
bigger in magnitude than all the others, then for large $k$ almost every vector,
when hit repeatedly by $\mat{A}$, lines up with that dominant eigenvector. Write
any starting vector in the eigenbasis, $\vec{x} = \sum_i c_i\vec{v}_i$; then
\[
  \mat{A}^k\vec{x} = \sum_i c_i\,\lambda_i^k\,\vec{v}_i
  = \lambda_1^k\Big( c_1\vec{v}_1 + \sum_{i\ge 2} c_i\big(\tfrac{\lambda_i}{\lambda_1}\big)^k\vec{v}_i\Big).
\]
If $\abs{\lambda_1} > \abs{\lambda_i}$ for all $i\ge2$, every ratio
$(\lambda_i/\lambda_1)^k \to 0$, and the bracket collapses to $c_1\vec{v}_1$. The
process forgets where it started and aligns with $\vec{v}_1$
(Figure~\ref{ch24:fig-dominant}).

\begin{figure}[H]
\centering
\begin{tikzpicture}[scale=1.0,>=Stealth,font=\small]
  % the two eigen-axes
  \draw[->,vecblue,thick] (0,0) -- (3.2,0.0) node[right]{$\vec{v}_1$ (dominant)};
  \draw[->,vecred,thick]  (0,0) -- (0.0,2.6) node[above]{$\vec{v}_2$};
  % successive images of x: each step shrinks the v2 part, grows v1 part
  % start somewhat off-axis, end nearly on v1
  \coordinate (p0) at (0.5,2.0);
  \coordinate (p1) at (1.1,1.2);
  \coordinate (p2) at (1.9,0.55);
  \coordinate (p3) at (2.7,0.18);
  \foreach \p/\lab in {p0/$\vec{x}$,p1/$\mat{A}\vec{x}$,p2/$\mat{A}^2\vec{x}$,p3/$\mat{A}^3\vec{x}$}{
    \draw[->,structurepurple,thick] (0,0) -- (\p);
  }
  \node[above right] at (p0) {$\vec{x}$};
  \node[above right] at (p1) {$\mat{A}\vec{x}$};
  \node[above right] at (p2) {$\mat{A}^2\vec{x}$};
  \node[below right] at (p3) {$\mat{A}^3\vec{x}$};
  \draw[gray!55,dashed] (p0) to[bend left=12] (p3);
\end{tikzpicture}
\caption{Dominant-eigenvalue behavior. Repeatedly applying $\mat{A}$ shrinks the
$\vec{v}_2$ component (by the factor $\lambda_2/\lambda_1$ each step) while the
$\vec{v}_1$ component grows fastest, so the images swing toward the dominant
eigen-direction $\vec{v}_1$. This is the engine behind PageRank, the power
method, and Markov-chain convergence.}
\label{ch24:fig-dominant}
\end{figure}

\begin{applicationbox}
\textbf{Google's PageRank and Markov chains.} A web-surfing model is a giant
matrix $\mat{A}$ whose largest eigenvalue is $\lambda_1 = 1$, with all other
$\abs{\lambda_i} < 1$. The formula above says $\mat{A}^k\vec{x}$ converges, for
\emph{any} starting distribution $\vec{x}$, to a multiple of the dominant
eigenvector $\vec{v}_1$ -- the \textbf{stationary distribution}. That eigenvector
\emph{is} the PageRank vector: the long-run fraction of time a random surfer
spends on each page. Diagonalization is not just a computational shortcut here;
it is the explanation of \emph{why} the ranking exists and is unique. The
``second eigenvalue'' $\abs{\lambda_2}$ controls how fast the algorithm converges
-- smaller $\abs{\lambda_2}$ means PageRank settles in fewer iterations. The same
mathematics governs how quickly a shuffled deck becomes random and how fast an
MCMC sampler in modern Bayesian statistics mixes.
\end{applicationbox}

\section{The Geometry: Three Clean Steps}
\label{ch24:geometry}

The formula $\mat{A} = \mat{S}\mat{\Lambda}\mat{S}\inv$ is best understood as a
\textbf{sentence read right to left}, because matrices act on a vector from the
right. To compute $\mat{A}\vec{x}$ you first apply $\mat{S}\inv$, then
$\mat{\Lambda}$, then $\mat{S}$:
\[
  \mat{A}\vec{x}
  = \mat{S}\,\big(\,\mat{\Lambda}\,(\,\mat{S}\inv\vec{x}\,)\,\big).
\]
Each factor has a clean geometric job.

\begin{itemize}
  \item $\mat{S}\inv$: \textbf{change into the eigenbasis.} It rewrites $\vec{x}$
        in terms of how much of each eigenvector it contains -- it reports the
        coordinates $(c_1,\dots,c_n)$ where $\vec{x}=\sum c_i\vec{v}_i$.
  \item $\mat{\Lambda}$: \textbf{stretch each eigen-axis.} In those coordinates,
        scale the $i$-th by $\lambda_i$. Simple, decoupled, no mixing.
  \item $\mat{S}$: \textbf{change back to standard coordinates.} It rebuilds the
        ordinary vector from its (now-stretched) eigen-coordinates.
\end{itemize}

\begin{figure}[H]
\centering
\begin{tikzpicture}[scale=0.92,>=Stealth,font=\small]
  % stage boxes
  \tikzstyle{stg}=[draw,rounded corners,minimum width=2.3cm,minimum height=1.3cm,align=center]
  \node[stg] (a) at (0,0) {standard\\$\vec{x}$};
  \node[stg] (b) at (3.7,0) {eigen-coords\\$\mat{S}\inv\vec{x}$};
  \node[stg] (c) at (7.6,0) {stretched\\$\mat{\Lambda}\mat{S}\inv\vec{x}$};
  \node[stg] (d) at (11.5,0) {standard\\$\mat{A}\vec{x}$};
  \draw[->,thick,vecblue] (a) -- node[above]{$\mat{S}\inv$} node[below]{\scriptsize translate in} (b);
  \draw[->,thick,vecred]  (b) -- node[above]{$\mat{\Lambda}$} node[below]{\scriptsize scale axes} (c);
  \draw[->,thick,vecblue] (c) -- node[above]{$\mat{S}$} node[below]{\scriptsize translate out} (d);
\end{tikzpicture}
\caption{Reading $\mat{A}=\mat{S}\mat{\Lambda}\mat{S}\inv$ as a pipeline:
translate into the eigen-language, do the easy stretching there, translate back.
Every hard thing $\mat{A}$ does is just easy stretching wearing a disguise.}
\label{ch24:fig-pipeline}
\end{figure}

This is exactly the change-of-basis story you will meet again in the chapter on
linear transformations: $\mat{\Lambda} = \mat{S}\inv\mat{A}\mat{S}$ says
\emph{the same machine $\mat{A}$, described in the eigenvector coordinate system,
is the diagonal matrix $\mat{\Lambda}$}. The matrix and its diagonal form are the
same transformation seen through two different windows.

\begin{connectionbox}
\textbf{Similar matrices and ``the same map in new clothes.''} Two matrices
$\mat{A}$ and $\mat{B}$ are called \textbf{similar} if $\mat{B} =
\mat{P}\inv\mat{A}\mat{P}$ for some invertible $\mat{P}$. Similarity is exactly
``same linear map, different basis.'' Diagonalization is the special, lucky case
where one of the two similar matrices happens to be diagonal. Similar matrices
\emph{share} their eigenvalues, their determinant ($\det\mat{A} = \prod\lambda_i$),
their trace ($\tr\mat{A} = \sum\lambda_i$), and their rank -- these are
properties of the underlying transformation, not of the coordinate window you
view it through. Two chapters from now this idea blossoms into the Jordan form,
which asks: when a matrix \emph{cannot} be made diagonal, what is the simplest
form it \emph{can} be made similar to?
\end{connectionbox}

\begin{walkbox}
Take a walk and chew on this: the determinant of $\mat{A}$ equals the product of
its eigenvalues, and the trace equals their sum. From $\mat{A} =
\mat{S}\mat{\Lambda}\mat{S}\inv$, can you see \emph{why}? Hint: $\det$ turns
products into products and $\det(\mat{S}\inv)=1/\det(\mat{S})$, so the $\mat{S}$'s
cancel; for the trace, recall the cyclic property $\tr(\mat{X}\mat{Y}) =
\tr(\mat{Y}\mat{X})$. Both invariants are reading the diagonal of
$\mat{\Lambda}$.
\end{walkbox}

\section{Which Matrices Diagonalize? Independent Eigenvectors}
\label{ch24:which}

Everything above needed one thing: $\mat{S}$ must be invertible, i.e. the
eigenvectors must be \emph{linearly independent}. When does that happen? The key
fact does most of the work for us.

\begin{theorem}[Distinct eigenvalues give independent eigenvectors]
\label{ch24:thm-distinct}
Eigenvectors corresponding to \emph{distinct} eigenvalues are linearly
independent. In particular, if an $n\times n$ matrix has $n$ distinct
eigenvalues, it has $n$ independent eigenvectors and is therefore
diagonalizable.
\end{theorem}

\begin{proof}[Sketch for two eigenvalues]
Suppose $\vec{v}_1,\vec{v}_2$ have distinct eigenvalues $\lambda_1\neq\lambda_2$
and some combination vanishes: $c_1\vec{v}_1 + c_2\vec{v}_2 = \vec0$. Apply
$\mat{A}$ to get $c_1\lambda_1\vec{v}_1 + c_2\lambda_2\vec{v}_2 = \vec0$. Now
subtract $\lambda_1$ times the first equation: the $\vec{v}_1$ term cancels,
leaving $c_2(\lambda_2 - \lambda_1)\vec{v}_2 = \vec0$. Since
$\lambda_2-\lambda_1\neq0$ and $\vec{v}_2\neq\vec0$, we must have $c_2 = 0$, and
then the first equation forces $c_1 = 0$. No nontrivial combination vanishes, so
the vectors are independent. The same ``apply $\mat{A}$ and subtract'' argument
extends to any number of distinct eigenvalues by induction.
\end{proof}

So \textbf{distinct eigenvalues guarantee diagonalizability}. This is the common,
comfortable case. But the converse is false, and reading it correctly will save
you from two classic traps.

\begin{trapbox}
\textbf{Trap 1: repeated eigenvalues do not doom a matrix.} You can have repeated
eigenvalues and still be perfectly diagonalizable. The identity matrix $\mat{I}$
has eigenvalue $1$ repeated $n$ times, yet \emph{every} vector is an eigenvector
-- it is already diagonal! A repeated eigenvalue is only a \emph{warning} that
you must go check whether you can find enough independent eigenvectors for it.

\textbf{Trap 2: diagonalizability is about eigenVECTORS, not eigenVALUES.} The
question is never ``are there $n$ eigenvalues?'' (counted with repetition there
always are, over $\C$). The question is ``are there $n$ independent
eigenVECTORS?'' Those are different questions, and the gap between them is the
whole subject of the next section.
\end{trapbox}

\subsection{Algebraic versus geometric multiplicity}
\label{ch24:multiplicity}

To say precisely when a repeated eigenvalue is harmless, we need two counts.

\begin{definition}[Two multiplicities]
\label{ch24:def-mult}
For an eigenvalue $\lambda$ of $\mat{A}$:
\begin{itemize}
  \item its \textbf{algebraic multiplicity} is how many times $\lambda$ appears
        as a root of the characteristic polynomial $\det(\mat{A}-\lambda\mat{I})$;
  \item its \textbf{geometric multiplicity} is the dimension of its eigenspace
        $\nullsp(\mat{A}-\lambda\mat{I})$ -- the number of independent
        eigenvectors you can actually find for $\lambda$.
\end{itemize}
Always $1 \le (\text{geometric}) \le (\text{algebraic})$.
\end{definition}

The matrix is diagonalizable exactly when there is no shortfall anywhere -- when
geometric multiplicity equals algebraic multiplicity for every eigenvalue, so the
eigenvectors add up to a full basis of $n$.

\begin{formalbox}
\textbf{Diagonalizability criterion.} $\mat{A}$ is diagonalizable (over a field
containing all its eigenvalues) if and only if, for every eigenvalue $\lambda$,
\[
  \dim\nullsp(\mat{A}-\lambda\mat{I}) \;=\; \text{(algebraic multiplicity of }\lambda).
\]
Equivalently, the geometric multiplicities sum to $n$. When some eigenvalue has
geometric multiplicity strictly less than its algebraic multiplicity, the
eigenvectors fall short of a basis and the matrix is \textbf{defective} -- it
cannot be diagonalized.
\end{formalbox}

\section{Defective Matrices: When Diagonalization Fails}
\label{ch24:defective}

Some matrices simply do not have enough eigenvectors, and no cleverness can
manufacture them. The smallest and most important example is the shear.

\begin{example}[The shear that cannot be diagonalized]
\label{ch24:shear}
Consider
\[
  \mat{A} = \begin{bmatrix} 1 & 1 \\ 0 & 1 \end{bmatrix}.
\]
It is upper triangular, so the eigenvalues are on the diagonal: $\lambda = 1$,
twice. Algebraic multiplicity of $\lambda=1$ is $2$. Now hunt for eigenvectors by
solving $(\mat{A}-\mat{I})\vec{v} = \vec0$:
\[
  \mat{A} - \mat{I} = \begin{bmatrix} 0 & 1 \\ 0 & 0 \end{bmatrix},
  \qquad
  \begin{bmatrix} 0 & 1 \\ 0 & 0 \end{bmatrix}\colvec{v_1\\v_2} = \vec0
  \;\Longrightarrow\; v_2 = 0.
\]
The only solutions are multiples of $\colvec{1\\0}$. The eigenspace is a single
line -- geometric multiplicity $1$, but algebraic multiplicity $2$. There is a
shortfall of one. We can never assemble two independent eigenvectors, so
$\mat{S}$ cannot be formed, and $\mat{A}$ is \textbf{defective}: not
diagonalizable.
\end{example}

Why does this matrix resist? Geometrically a shear slides the plane sideways: the
horizontal axis is pinned (that is the one eigenvector), but \emph{no other
direction is merely scaled} -- every other vector gets tilted toward the
horizontal. There is no second axis that the machine leaves pointing the same
way, so there is no second eigenvector, and hence no eigenbasis.

\begin{figure}[H]
\centering
\begin{tikzpicture}[scale=1.0,>=Stealth,font=\small]
  \draw[->,gray!55] (-0.3,0) -- (4.0,0) node[right]{$x$};
  \draw[->,gray!55] (0,-0.3) -- (0,2.6) node[above]{$y$};
  % the pinned eigen-direction (x-axis): a vector and its image coincide in direction
  \draw[->,vecblue,very thick] (0,0) -- (2.4,0) node[below=2pt]{$\vec{v}=\colvec{1\\0}$ pinned};
  % a generic vector and its sheared image, tilting toward horizontal
  \draw[->,vecred,very thick] (0,0) -- (1.0,1.5) node[above left]{$\vec{u}$};
  \draw[->,vecred!55,very thick,dashed] (0,0) -- (2.5,1.5) node[right]{$\mat{A}\vec{u}$ (tilted, not scaled)};
  \draw[gray!50,dashed] (1.0,1.5) -- (2.5,1.5);
  \node[align=center] at (2.6,2.3) {\scriptsize every non-horizontal\\[-2pt]\scriptsize arrow tips over};
\end{tikzpicture}
\caption{The shear $\protect\begin{bmatrix}1&1\\0&1\protect\end{bmatrix}$ pins
the horizontal axis (its one eigenvector) but tilts every other vector toward the
horizontal -- changing its \emph{direction}, not merely its length. With only one
eigen-direction there is no eigenbasis, so the matrix is defective.}
\label{ch24:fig-shear}
\end{figure}

\begin{intuitionbox}
A defective matrix is one that does something a diagonal matrix simply cannot:
it has a direction that gets \emph{rotated into} another, never purely scaled.
Pure scaling along independent axes is the entire repertoire of a diagonal
matrix, so any genuine ``tipping over'' that cannot be undone by a change of axes
is the fingerprint of defectiveness. The repeated-eigenvalue shear is the
canonical culprit.
\end{intuitionbox}

\begin{trapbox}
A repeated eigenvalue \emph{may or may not} be defective -- you must check.
Compare these two, both with eigenvalue $5$ doubled:
\[
  \begin{bmatrix} 5 & 0 \\ 0 & 5 \end{bmatrix}
  \ \text{(geom.\ mult.\ }2\text{, diagonalizable),}
  \qquad
  \begin{bmatrix} 5 & 1 \\ 0 & 5 \end{bmatrix}
  \ \text{(geom.\ mult.\ }1\text{, defective).}
\]
Same characteristic polynomial $(\lambda-5)^2$, opposite fates. The off-diagonal
$1$ is what kills the second eigenvector. Never judge diagonalizability from the
eigenvalues alone.
\end{trapbox}

\begin{remark}[There is still a ``best possible'' form]
\label{ch24:jordan-preview}
A defective matrix cannot be diagonalized, but it is not the end of the road. The
\textbf{Jordan form} (a later chapter) is the nearest substitute: it is
\emph{almost} diagonal, with the eigenvalues on the diagonal and a few stray
$1$'s just above it -- exactly one $1$ for each missing eigenvector. The shear
$\begin{bsmallmatrix}1&1\\0&1\end{bsmallmatrix}$ is its own Jordan form: that
stray $1$ is the visible scar of the eigenvector it lacks.
\end{remark}

\begin{connectionbox}
\textbf{Defectiveness is fragile, and that is a story about calculus.} Defective
matrices are rare in the sense that an arbitrarily tiny perturbation breaks the
repeated eigenvalue apart into distinct ones, instantly making the matrix
diagonalizable. For instance $\begin{bsmallmatrix}1&1\\0&1\end{bsmallmatrix}$ is
defective, but $\begin{bsmallmatrix}1&1\\0&1+\eps\end{bsmallmatrix}$ has the two
distinct eigenvalues $1$ and $1+\eps$, hence two eigenvectors, for every
$\eps\neq0$. As $\eps\to0$ those two eigenvectors swing toward each other and
\emph{collide}, and $\mat{S}\inv$ blows up. This is why numerical software is
nervous near defective matrices: the eigenvectors become almost parallel and
$\mat{S}$ becomes nearly singular. The cure that everyone eventually reaches for
is the singular value decomposition, which never has this problem -- but that is
chapters away.
\end{connectionbox}

\section{Application: Solving the Fibonacci Recurrence}
\label{ch24:fibonacci}

Diagonalization turns a recurrence -- a rule that builds each term from the last
ones -- into a matrix power, which we now know how to compute exactly. The most
famous recurrence in the world is Fibonacci:
\[
  F_0 = 0,\quad F_1 = 1,\qquad F_{k+1} = F_k + F_{k-1}.
\]
So $0,1,1,2,3,5,8,13,21,\dots$. Is there a \emph{closed form} for $F_k$ -- a
formula that jumps straight to the thousandth term without computing the first
nine hundred ninety-nine?

\textbf{Step 1: make it a matrix.} The trick for any recurrence is to track a
little \emph{state vector} holding enough recent terms to take the next step.
Here, two terms suffice. Let $\vec{u}_k = \colvec{F_{k+1}\\ F_k}$. The recurrence
says
\[
  \colvec{F_{k+1}\\F_k} = \begin{bmatrix} 1 & 1 \\ 1 & 0 \end{bmatrix}\colvec{F_k\\F_{k-1}},
  \qquad\text{i.e.}\qquad
  \vec{u}_k = \mat{A}\,\vec{u}_{k-1},
  \quad
  \mat{A} = \begin{bmatrix} 1 & 1 \\ 1 & 0 \end{bmatrix}.
\]
Check the top row: $F_{k+1} = 1\cdot F_k + 1\cdot F_{k-1}$, the recurrence
itself. The bottom row just copies $F_k$ down so it is ready to be the ``previous
term'' next time. Iterating from $\vec{u}_0 = \colvec{F_1\\F_0} = \colvec{1\\0}$,
\[
  \vec{u}_k = \mat{A}^k\,\vec{u}_0.
\]
A closed form for $F_k$ is now a closed form for $\mat{A}^k$ -- and we know
exactly how to get that.

\textbf{Step 2: diagonalize $\mat{A}$.} Characteristic polynomial:
\[
  \det\begin{bmatrix} 1-\lambda & 1 \\ 1 & -\lambda \end{bmatrix}
  = (1-\lambda)(-\lambda) - 1
  = \lambda^2 - \lambda - 1 = 0.
\]
The roots are the golden ratio and its conjugate:
\[
  \lambda_1 = \varphi = \frac{1+\sqrt5}{2} \approx 1.618,
  \qquad
  \lambda_2 = \psi = \frac{1-\sqrt5}{2} \approx -0.618.
\]
They are distinct, so $\mat{A}$ is diagonalizable. For eigenvectors, from
$(\mat{A}-\lambda\mat{I})\vec{v}=\vec0$ the bottom row gives $v_1 = \lambda v_2$,
so an eigenvector for $\lambda$ is $\colvec{\lambda\\1}$. Thus
\[
  \mat{S} = \begin{bmatrix} \varphi & \psi \\ 1 & 1 \end{bmatrix},
  \qquad
  \mat{\Lambda} = \begin{bmatrix} \varphi & 0 \\ 0 & \psi \end{bmatrix}.
\]

\textbf{Step 3: extract the formula.} We want $F_k$, the bottom entry of
$\vec{u}_k = \mat{S}\mat{\Lambda}^k\mat{S}\inv\vec{u}_0$. Rather than grind the
full product, use the cleaner eigenbasis route: expand the start vector
$\vec{u}_0 = \colvec{1\\0}$ in eigenvectors, $\vec{u}_0 = c_1\colvec{\varphi\\1} +
c_2\colvec{\psi\\1}$. The two equations are $c_1\varphi + c_2\psi = 1$ and
$c_1 + c_2 = 0$, giving $c_2 = -c_1$ and $c_1(\varphi-\psi) = 1$. Since
$\varphi - \psi = \sqrt5$, we get $c_1 = 1/\sqrt5$ and $c_2 = -1/\sqrt5$. Then
\[
  \vec{u}_k = \mat{A}^k\vec{u}_0
  = c_1\varphi^k\colvec{\varphi\\1} + c_2\psi^k\colvec{\psi\\1}.
\]
Reading the bottom entry ($F_k$) gives the celebrated \textbf{Binet formula}:
\[
  \boxed{\;F_k = \frac{\varphi^k - \psi^k}{\sqrt5}
        = \frac{1}{\sqrt5}\!\left[\left(\frac{1+\sqrt5}{2}\right)^{\!k}
        - \left(\frac{1-\sqrt5}{2}\right)^{\!k}\right].\;}
\]

\begin{example}[Sanity check on small $k$]
\label{ch24:binet-check}
A formula full of irrational $\sqrt5$'s had better produce integers. Test
$k=5$. With $\varphi^5 \approx 11.0902$ and $\psi^5 \approx -0.0902$,
\[
  F_5 = \frac{11.0902 - (-0.0902)}{\sqrt5}
      = \frac{11.1803}{2.2360} = 5.\ \checkmark
\]
Indeed the sequence is $0,1,1,2,3,5,\dots$ and $F_5 = 5$. All the $\sqrt5$
messiness cancels to leave a whole number every time -- a small miracle that the
two conjugate eigenvalues conspire to produce.
\end{example}

\begin{intuitionbox}
Binet's formula \emph{is} the eigenbasis decomposition $\mat{A}^k\vec{u}_0 =
\sum_i c_i\lambda_i^k\vec{v}_i$ written out for one specific matrix. The Fibonacci
numbers are a blend of two pure exponential modes, $\varphi^k$ and $\psi^k$.
Since $\abs{\psi}\approx0.618<1$, the $\psi^k$ mode shrinks to nothing, so for
large $k$ the Fibonacci numbers are essentially $\varphi^k/\sqrt5$ -- they grow
by a factor of $\varphi$ each step (Figure~\ref{ch24:fig-fibmodes}). That is precisely the
``dominant-eigenvalue'' long-run behavior of
Section~\ref{ch24:longrun}, and it is why consecutive Fibonacci ratios
$F_{k+1}/F_k$ march toward the golden ratio.
\end{intuitionbox}

\begin{figure}[H]
\centering
\begin{tikzpicture}[>=Stealth,font=\small]
\begin{axis}[
  width=10.5cm, height=6cm,
  xlabel={$k$}, ylabel={value},
  xmin=0, xmax=8, ymin=-1.5, ymax=22,
  axis lines=left,
  legend pos=north west, legend cell align=left,
  grid=major, grid style={gray!18},
]
  % phi^k / sqrt5 mode
  \addplot[vecblue,thick,domain=0:8,samples=80]{(((1+sqrt(5))/2)^x)/sqrt(5)};
  \addlegendentry{$\varphi^k/\sqrt5$ (growing mode)}
  % psi^k / sqrt5 mode (tiny, alternating)
  \addplot[vecred,thick,domain=0:8,samples=80]{(((1-sqrt(5))/2)^x)/sqrt(5)};
  \addlegendentry{$\psi^k/\sqrt5$ (shrinking mode)}
  % actual Fibonacci integers
  \addplot[only marks,mark=*,mark size=1.6pt,structurepurple]
    coordinates {(0,0)(1,1)(2,1)(3,2)(4,3)(5,5)(6,8)(7,13)(8,21)};
  \addlegendentry{$F_k$ (integers)}
\end{axis}
\end{tikzpicture}
\caption{Fibonacci as two eigen-modes. $F_k = \varphi^k/\sqrt5 - \psi^k/\sqrt5$ is
the difference of a fast-growing mode $\varphi^k/\sqrt5$ (the dominant
eigenvalue) and a tiny alternating mode $\psi^k/\sqrt5$ that decays to nothing.
The integers $F_k$ hug the growing mode almost perfectly -- which is why $F_k$ is
the nearest integer to $\varphi^k/\sqrt5$.}
\label{ch24:fig-fibmodes}
\end{figure}

\begin{applicationbox}
\textbf{Every linear recurrence yields to this.} The Fibonacci move is completely
general. Any constant-coefficient linear recurrence
$a_{k+1} = p_1 a_k + p_2 a_{k-1} + \dots + p_d a_{k-d+1}$ becomes
$\vec{u}_k = \mat{A}\vec{u}_{k-1}$ with a $d\times d$ \emph{companion matrix}
$\mat{A}$, and diagonalizing $\mat{A}$ produces a closed form as a sum of
$\lambda_i^k$ modes. This is exactly how engineers analyze the stability of
discrete-time control systems and digital filters (stable iff every
$\abs{\lambda_i}<1$), how population biologists project age-structured
populations with Leslie matrices, and how the runtime of a divide-and-conquer
algorithm is read off from the roots of its recurrence. One factorization, a
whole field of consequences.
\end{applicationbox}

\section{The Recipe and the Cautions}
\label{ch24:recipe}

Pulling the chapter together, here is the full procedure and the places people
stumble.

\begin{focusbox}
\textbf{To diagonalize an $n\times n$ matrix $\mat{A}$:}
\begin{enumerate}
  \item Find all eigenvalues: solve $\det(\mat{A}-\lambda\mat{I}) = 0$.
  \item For each $\lambda$, find a basis of its eigenspace
        $\nullsp(\mat{A}-\lambda\mat{I})$.
  \item Count: if the eigenvectors total $n$ (every geometric multiplicity meets
        its algebraic multiplicity), you can diagonalize. If not, $\mat{A}$ is
        defective -- stop.
  \item Put the $n$ eigenvectors as columns of $\mat{S}$, their eigenvalues
        (same order!) on the diagonal of $\mat{\Lambda}$. Then $\mat{A} =
        \mat{S}\mat{\Lambda}\mat{S}\inv$, and $\mat{A}^k =
        \mat{S}\mat{\Lambda}^k\mat{S}\inv$.
\end{enumerate}
\end{focusbox}

\begin{trapbox}
\textbf{Four diagonalization traps, collected.}
\begin{itemize}
  \item \textbf{Order mismatch.} The $j$-th column of $\mat{S}$ and the $j$-th
        diagonal entry of $\mat{\Lambda}$ must be a matched eigenpair. Shuffle
        one without the other and $\mat{S}\mat{\Lambda}\mat{S}\inv \neq \mat{A}$.
  \item \textbf{Confusing the two multiplicities.} ``Three eigenvalues'' is not
        the test; ``three independent eigenvectors'' is.
  \item \textbf{Forgetting complex eigenvalues.} A rotation like
        $\begin{bsmallmatrix}0&-1\\1&0\end{bsmallmatrix}$ has \emph{no real}
        eigenvectors (it turns every real vector), so it is not diagonalizable
        over $\R$ -- but it \emph{is} diagonalizable over $\C$, with eigenvalues
        $\pm i$. Always note which field you are working over.
  \item \textbf{$\mat{\Lambda}^k$ shortcut applies only to diagonal
        $\mat{\Lambda}$.} You may raise diagonal entries to the $k$ individually
        \emph{only because} $\mat{\Lambda}$ is diagonal. You cannot do this to
        $\mat{A}$ directly.
\end{itemize}
\end{trapbox}

\begin{strangbox}
After Strang, the natural next questions are: which matrices are
\emph{guaranteed} a full set of eigenvectors -- indeed orthogonal ones -- so that
$\mat{S}\inv$ is just $\mat{S}\T$? That is the \textbf{spectral theorem} for
symmetric matrices, the cleanest diagonalization of all, coming in a later
chapter. And what replaces $\mat{S}\mat{\Lambda}\mat{S}\inv$ for rectangular or
defective matrices, where eigenvectors fail us? That is the \textbf{singular
value decomposition}, the crown jewel of the second half of this book.
Diagonalization is the doorway to both.
\end{strangbox}

\subsection{Cost and intuition recap}
\label{ch24:recap}

\begin{recoverybox}
\textbf{The one-line intuition.} Eigenvectors are the axes in which $\mat{A}$ is
diagonal; $\mat{A} = \mat{S}\mat{\Lambda}\mat{S}\inv$ says ``translate into those
axes, scale, translate back.''

\textbf{The cost.} Diagonalizing once is real work: forming the characteristic
polynomial and solving for eigenvalues is the expensive step (for an $n\times n$
matrix, a full eigendecomposition costs on the order of $n^3$ arithmetic
operations, the same order as one matrix inversion). But that cost is paid
\emph{once}. After it,
\begin{itemize}
  \item a single power $\mat{A}^k$ costs only $n$ scalar powers $\lambda_i^k$ plus
        two fixed matrix products to sandwich them -- the price no longer grows
        with $k$, whereas naive repeated multiplication costs about $k$ matrix
        products (order $k\,n^3$);
  \item the inverse is free: $\mat{A}\inv =
        \mat{S}\mat{\Lambda}\inv\mat{S}\inv$ just reciprocates the eigenvalues;
  \item the determinant ($\prod_i\lambda_i$) and trace ($\sum_i\lambda_i$) are
        read straight off $\mat{\Lambda}$.
\end{itemize}

\textbf{The catch.} You need $n$ independent eigenvectors. Distinct eigenvalues
guarantee them; repeated eigenvalues require a check (geometric multiplicity must
meet algebraic multiplicity), and a defective matrix has no eigenbasis at all.
\end{recoverybox}

\section*{Exercises}
\addcontentsline{toc}{section}{Exercises}

\begin{exercise}[Warm-up factorization]
\label{ch24:ex-warmup}
Let $\mat{A} = \begin{bsmallmatrix} 4 & 0 \\ 1 & 2 \end{bsmallmatrix}$. Find its
eigenvalues and eigenvectors, write down $\mat{S}$ and $\mat{\Lambda}$, and verify
by direct multiplication that $\mat{S}\mat{\Lambda}\mat{S}\inv = \mat{A}$.
\end{exercise}

\begin{exercise}[A power]
\label{ch24:ex-power}
Using the diagonalization of $\mat{A} = \begin{bsmallmatrix} 2 & 1 \\ 1 & 2
\end{bsmallmatrix}$ from Section~\ref{ch24:first-example}, compute $\mat{A}^5$
exactly via $\mat{S}\mat{\Lambda}^5\mat{S}\inv$. (Check: the diagonal entries
should be $\tfrac{1+3^5}{2} = 122$.)
\end{exercise}

\begin{exercise}[Defective or not?]
\label{ch24:ex-defective}
For each matrix, decide whether it is diagonalizable over $\R$, giving the
algebraic and geometric multiplicity of every eigenvalue:
\[
  \begin{bmatrix} 3 & 1 \\ 0 & 3 \end{bmatrix},
  \qquad
  \begin{bmatrix} 3 & 0 \\ 0 & 3 \end{bmatrix},
  \qquad
  \begin{bmatrix} 2 & 0 \\ 0 & 5 \end{bmatrix},
  \qquad
  \begin{bmatrix} 0 & -1 \\ 1 & 0 \end{bmatrix}.
\]
\end{exercise}

\begin{exercise}[Trace and determinant from $\mat{\Lambda}$]
\label{ch24:ex-tracedet}
Using $\mat{A} = \mat{S}\mat{\Lambda}\mat{S}\inv$ together with
$\tr(\mat{X}\mat{Y}) = \tr(\mat{Y}\mat{X})$ and $\det(\mat{X}\mat{Y}) =
\det(\mat{X})\det(\mat{Y})$, prove carefully that $\tr\mat{A} = \sum_i\lambda_i$
and $\det\mat{A} = \prod_i\lambda_i$ for any diagonalizable $\mat{A}$.
\end{exercise}

\begin{exercise}[A recurrence of your own]
\label{ch24:ex-recurrence}
Define $a_0 = 1$, $a_1 = 1$, and $a_{k+1} = 5a_k - 6a_{k-1}$. Write the state
recurrence $\vec{u}_k = \mat{A}\vec{u}_{k-1}$ with $\vec{u}_k =
\colvec{a_{k+1}\\a_k}$, diagonalize the $2\times2$ matrix $\mat{A}$ (its
eigenvalues are nice integers), and derive a closed form $a_k = \alpha\,2^k +
\beta\,3^k$. Determine $\alpha,\beta$ from the initial conditions.
\end{exercise}

\begin{exercise}[Inverse from the factorization]
\label{ch24:ex-inverse}
Show that if $\mat{A} = \mat{S}\mat{\Lambda}\mat{S}\inv$ is invertible (all
$\lambda_i\neq0$), then $\mat{A}\inv = \mat{S}\mat{\Lambda}\inv\mat{S}\inv$ with
$\mat{\Lambda}\inv = \diag(1/\lambda_1,\dots,1/\lambda_n)$. Conclude that
$\mat{A}$ and $\mat{A}\inv$ share the same eigenvectors, with reciprocal
eigenvalues.
\end{exercise}

\begin{exercise}[Commuting matrices, shared eigenvectors]
\label{ch24:ex-commute}
Suppose $\mat{A}$ is diagonalizable with \emph{distinct} eigenvalues and
$\mat{A}\mat{B} = \mat{B}\mat{A}$. Show that every eigenvector of $\mat{A}$ is
also an eigenvector of $\mat{B}$. (Hint: if $\mat{A}\vec{v} = \lambda\vec{v}$,
apply $\mat{A}$ to $\mat{B}\vec{v}$ and use commuting to show $\mat{B}\vec{v}$
lies in the same eigenspace, which is one-dimensional.) This is why commuting
matrices can be ``simultaneously diagonalized.''
\end{exercise}

\begin{exercise}[Building a matrix from its spectrum]
\label{ch24:ex-build}
Construct a $2\times2$ matrix whose eigenvectors are $\colvec{2\\1}$ (eigenvalue
$4$) and $\colvec{1\\1}$ (eigenvalue $-1$), by forming
$\mat{S}\mat{\Lambda}\mat{S}\inv$ explicitly. Then verify your matrix really has
those eigenpairs.
\end{exercise}

\begin{exercise}[Why $\psi^k$ is invisible]
\label{ch24:ex-fib-asymp}
From Binet's formula, prove that $F_k$ is the \emph{nearest integer} to
$\varphi^k/\sqrt5$ for all $k\ge0$. (Hint: bound $\abs{\psi^k}/\sqrt5 <
\tfrac12$.) Use this to argue that $F_{k+1}/F_k \to \varphi$, and connect this to
the dominant-eigenvalue discussion of Section~\ref{ch24:longrun}.
\end{exercise}

\begin{exercise}[A $3\times3$ challenge]
\label{ch24:ex-3x3}
Diagonalize
\[
  \mat{A} = \begin{bmatrix} 2 & 0 & 0 \\ 1 & 2 & 0 \\ 0 & 0 & 3 \end{bmatrix}.
\]
Find all eigenvalues and a basis of each eigenspace. Is $\mat{A}$ diagonalizable?
Identify the algebraic and geometric multiplicities of $\lambda = 2$, and explain
in one sentence what the off-diagonal $1$ does to your conclusion.
\end{exercise}
